\chapter{Methodology}

\section{Overview of Technical Approach}
\label{sec:method_overview}

TechPulse has been implemented as a cloud-native Hybrid Retrieval-Augmented Generation system designed to synthesize emerging technology intelligence from heterogeneous sources. The technical approach is organized around five interconnected components: (1) continuous heterogeneous data ingestion from five sources; (2) deterministic preprocessing and semantic embedding; (3) hybrid multi-stage retrieval; (4) context-bound grounded generation with automatic LLM fallback; and (5) production-grade MLOps deployment and monitoring.

The core hypotheses under evaluation are stated formally as:

\begin{itemize}[leftmargin=0.5in]
  \item \textbf{H\textsubscript{0}}: There is no statistically significant difference in mean context precision between hybrid and baseline (vector-only) retrieval, evaluated over 50 matched queries under identical corpus and embedding conditions.
  \item \textbf{H\textsubscript{1}}: Hybrid retrieval achieves statistically higher mean context precision than baseline retrieval (one-sided Wilcoxon signed-rank test, $\alpha = 0.05$).
\end{itemize}

A secondary performance hypothesis requires that the hybrid retrieval pipeline does not exceed the p95 latency bound ($\tau = 2.0$~seconds for the retrieval stage) or the monthly cost ceiling (USD~15 at an assumed workload of 1,500 queries/month, i.e., USD~0.01 per query). The evaluation is designed to be informative regardless of outcome: if H\textsubscript{0} cannot be rejected, the study characterises the conditions under which hybrid and baseline converge---an empirically valuable finding given the widespread adoption of hybrid retrieval in production RAG systems without controlled comparative evidence. These hypotheses are tested through a controlled experiment comparing both configurations under identical corpus and embedding conditions.

\section{Why This Approach is Suitable}
\label{sec:method_rationale}

Three considerations motivate the hybrid RAG approach for this specific problem domain.

\textbf{Domain characteristics:} Emerging technology knowledge is temporally dynamic and distributed across heterogeneous sources. Vector-only retrieval treats all documents equally regardless of recency or source authority; hybrid retrieval with BM25 keyword search and RRF fusion explicitly addresses temporal drift, which is critical in rapidly evolving AI/ML domains where a paper from 18 months ago may already be superseded.

\textbf{Operational constraints:} The system must operate within cost and latency bounds suitable for a research prototype deployed on the AWS Free Tier. Full-corpus BM25 scoring adds $O(N)$ computation (one pass over $N$ pre-tokenized chunks) but is executed in-process against a cached index, avoiding network round-trips. Reciprocal Rank Fusion merges three ranked lists in $O(k)$, where $k$ is the candidate pool size (typically $4 \times \text{top\_k} = 20$). The combined overhead is negligible compared to the $O(Nd)$ baseline vector search and avoids the latency penalties of cross-encoder rerankers, which require a forward pass through a BERT-class model per candidate.

\textbf{Production feasibility:} The architecture uses AWS-managed services---Lambda, RDS PostgreSQL, API Gateway, EventBridge---eliminating the need to manage dedicated compute clusters or self-hosted vector databases. This enables a student team to deploy, iterate, and monitor a production-grade system within the course timeline and a prototype budget envelope.

\section{Technical Contributions and Design Decisions}
\label{sec:innovation}

The engineering contribution of TechPulse lies not in any single component but in their integration into a coherent, cost-governed MLOps system. Specifically, four design decisions represent well-motivated engineering choices that distinguish this system from typical RAG prototypes.

\textbf{True parallel hybrid retrieval with parameter-free fusion:} The system performs two independent retrieval passes---cosine vector search (pgvector) and BM25 keyword search (BM25Okapi over the full indexed corpus)---and fuses their rankings with Reciprocal Rank Fusion (RRF; Cormack et al., 2009). A third recency signal, computed as document age in days, is integrated as an additional RRF ranking. Because RRF operates on ordinal ranks rather than raw scores, it is parameter-free and scale-invariant: no weight tuning ($\alpha$/$\beta$/$\gamma$) is required, eliminating a common source of overfitting in hybrid retrieval systems. This parameter-free approach replaced an earlier design-phase weighted reranking formula ($\alpha$/$\beta$/$\gamma$) that proved sensitive to corpus composition and incurred prohibitive tail latency (p95~=~25\,s); the empirical evidence motivating this transition is documented in Section~\ref{sec:sensitivity}. The BM25 index is built over the full corpus so that IDF statistics are globally accurate, and refreshed via a 30-minute TTL cache to capture newly ingested documents without rebuilding on every query.

\textbf{Unified vector--relational retrieval without dedicated infrastructure:} By extending PostgreSQL with pgvector, TechPulse achieves hybrid retrieval (structured SQL metadata filtering + cosine vector search) within a single managed database, eliminating the operational complexity of running a separate vector store. This design choice reduces architectural surface area, simplifies deployment automation, and enables cost-governed scaling through standard RDS instances.

\textbf{Idempotent ingestion pipeline with SHA-256 content-addressed deduplication:} The SHA-256 hashing strategy applied to canonicalized document fields ensures that the pipeline can be re-run safely at any time without creating duplicate records. This often-overlooked but practically important engineering practice enables safe experiment replay, crash recovery, and corpus versioning without requiring a separate deduplication service.

\textbf{Budget-aware inference with programmatic halt:} The integration of AWS Cost Explorer queries into the inference path---where invocations are halted programmatically when monthly cost thresholds are approached---represents an operational optimization strategy not typically addressed in RAG research prototypes. This transforms cost from a passive monitoring concern into an active constraint enforced at runtime, consistent with production-grade MLOps principles.

\section{Baseline Retrieval Configuration}
\label{sec:baseline}

The baseline configuration performs dense cosine similarity search across the full indexed corpus. Given a query embedding $\mathbf{e}_q$ and document embeddings $\mathbf{e}_i$, cosine similarity is computed as:

\begin{equation}
  \text{Sim}(q, i) = \frac{\mathbf{e}_q \cdot \mathbf{e}_i}{\|\mathbf{e}_q\|\|\mathbf{e}_i\|}
  \label{eq:cosine}
\end{equation}

Top-$N$ documents are selected purely on similarity ranking with a minimum similarity threshold of 0.15 to filter retrieval noise. This low threshold is intentionally permissive in the baseline to avoid artificially handicapping it; borderline documents that would be excluded by a stricter threshold are instead ranked low by subsequent RRF fusion in the hybrid configuration. Under exact search, computational complexity scales as $O(Nd)$, where $N$ is the number of indexed chunks and $d = 384$ (the MiniLM embedding dimension). This configuration serves as the experimental control for retrieval performance comparison.

\section{Implemented Hybrid Retrieval Configuration}
\label{sec:hybrid}

The hybrid configuration introduces four refinement stages applied before generation.

\textbf{Stage 1 --- Metadata Filtering:} Structured SQL constraints reduce the candidate space based on document recency (a 180-day sliding window), source type, and document state. Only documents in the \texttt{INDEXED} state are considered. The 180-day window was chosen to balance corpus recency with sufficient document volume: at the expected ingestion rates across five sources, the window is projected to yield 500--2,000 indexed chunks, providing a meaningful retrieval pool while excluding content more than six months old, which is considered stale for fast-moving AI/ML topics. This pre-filtering stage reduces retrieval noise prior to semantic ranking.

\textit{Robustness to missing and noisy metadata.} The filter degrades gracefully when metadata is incomplete. The recency filter uses an inclusive \texttt{OR} clause (\texttt{published\_at IS NULL OR published\_at >= threshold}), retaining documents with missing dates rather than silently dropping them. Each source adapter normalizes metadata at ingestion (whitespace normalization, URL defaults, timestamp coercion from API-specific fields). In the RRF fusion stage, documents with \texttt{NULL} dates receive age $= 9{,}999$ days, deprioritizing them by recency while preserving their semantic and BM25 retrieval paths.

\textbf{Stage 2 --- Vector Similarity Search:} Cosine similarity (Equation~\ref{eq:cosine}) is applied over the filtered candidate subset using the pgvector \texttt{<=>} operator. To provide sufficient candidates for RRF fusion, $4 \times \text{top\_k}$ results are retrieved at this stage. The 4$\times$ multiplier ensures that the vector candidate list is large enough to overlap meaningfully with BM25 results during fusion, improving the diversity of the final merged ranking without inflating retrieval latency.

\textbf{Stage 3 --- Parallel BM25 Retrieval:} Concurrently with vector search, a full-corpus BM25 keyword search is executed using the BM25Okapi algorithm (Robertson \& Walker, 1994). Given a query $Q$ containing terms $q_1, q_2, \ldots, q_n$, the BM25 score for document $D$ is:

\begin{equation}
  \text{BM25}(D, Q) = \sum_{i=1}^{n} \text{IDF}(q_i) \cdot \frac{f(q_i, D) \cdot (k_1 + 1)}{f(q_i, D) + k_1 \cdot \bigl(1 - b + b \cdot \frac{|D|}{\text{avgdl}}\bigr)}
  \label{eq:bm25}
\end{equation}

\noindent where $f(q_i, D)$ is the term frequency of $q_i$ in $D$, $|D|$ is the document length, $\text{avgdl}$ is the average document length across the corpus, and $k_1 = 1.5$, $b = 0.75$ are the BM25Okapi defaults. Crucially, IDF is computed over the \emph{full} indexed corpus ($N = 6{,}220$ chunks at evaluation time), not over the filtered candidate set---computing IDF over a small candidate pool would produce up to $6\times$ error in term weighting (observed empirically during development). The BM25 index is built once and cached at the module level with a 30-minute TTL, refreshed automatically to capture newly ingested documents.

The same source-type and recency filters applied in Stage~1 are enforced on BM25 candidates in Python, ensuring both retrieval signals operate over an identical filtered corpus.

\textbf{Stage 4 --- Reciprocal Rank Fusion (RRF):} The union of vector and BM25 candidate sets is fused using Reciprocal Rank Fusion (Cormack, Clarke, \& Butt, 2009):

\begin{equation}
  \text{RRF}(d) = \sum_{r \in \mathcal{R}} \frac{1}{K + \text{rank}_r(d)}
  \label{eq:rrf}
\end{equation}

\noindent where $\mathcal{R} = \{\text{vector}, \text{BM25}, \text{recency}\}$ represents three independent rankings, $\text{rank}_r(d)$ is the 1-indexed rank of document $d$ in ranking $r$, and $K = 60$ is the standard RRF constant from the original paper. Recency is incorporated as a third ranking signal by sorting candidates by document age in days (ascending; documents with unknown publication dates receive age $= 9{,}999$ and sort last).

RRF has two key advantages over weighted linear scoring: (1)~it is \emph{parameter-free}---no $\alpha$/$\beta$/$\gamma$ weights require tuning, eliminating a common source of overfitting; and (2)~it is \emph{scale-invariant}---because it operates on ordinal ranks, it handles the incomparable score distributions of cosine similarity ($[0,1]$) and BM25 ($[0, \infty)$) without normalisation.

Figure~\ref{fig:retrieval_comparison} illustrates the design evolution across three retrieval strategies: the baseline vector-only pipeline (left), the original weighted hybrid that was superseded due to p95~=~25\,s tail latency and dataset-dependent weights (centre, faded), and the final BM25+RRF hybrid pipeline adopted in the production system (right).

% ---- TikZ: Three-way retrieval comparison diagram ----
\begin{figure}[H]
\centering
\resizebox{\textwidth}{!}{%
\begin{tikzpicture}[node distance=0.45cm and 1.4cm]

% ---- LEFT: BASELINE ----
\node[pbox=awsgray, minimum width=2.4cm] (bq) {User Query};
\node[pbox=awsblue, below=of bq, minimum width=2.4cm] (bvs) {Vector Similarity\\Search (Full Corpus)};
\node[pbox=awsgreen, below=of bvs, minimum width=2.4cm] (btopk) {Top-$k$ Selection};
\node[pbox=awsorange, below=of btopk, minimum width=2.4cm] (bgen) {LLM\\Generation};

\draw[parrow] (bq) -- (bvs);
\draw[parrow] (bvs) -- (btopk);
\draw[parrow] (btopk) -- (bgen);

\node[layerlabel=awsgray, above=0.3cm of bq, font=\small\bfseries] {\textbf{Baseline (Vector-Only)}};

\begin{scope}[on background layer]
\node[boundbox=awsgray, fit=(bq)(bvs)(btopk)(bgen), inner sep=8pt] {};
\end{scope}

% ---- CENTER: ORIGINAL HYBRID (Superseded) ----
\node[pbox=propgray, minimum width=2.4cm, right=4.0cm of bq] (oq) {User Query};
\node[pbox=propgray, below=of oq, minimum width=2.4cm] (omf) {Metadata Filter\\(recency + source)};
\node[pbox=propgray, below=of omf, minimum width=2.4cm] (ovs) {Vector Similarity\\Search (Filtered)};
\node[pbox=propred, below=of ovs, minimum width=2.4cm] (orr) {Weighted\\$\alpha$/$\beta$/$\gamma$ Reranking};
\node[pbox=propgray, below=of orr, minimum width=2.4cm] (otopk) {Top-$k$ Selection};
\node[pbox=propgray, below=of otopk, minimum width=2.4cm] (ogen) {LLM\\Generation};

\draw[parrow, color=propgray!60] (oq) -- (omf);
\draw[parrow, color=propgray!60] (omf) -- (ovs);
\draw[parrow, color=propgray!60] (ovs) -- (orr);
\draw[parrow, color=propgray!60] (orr) -- (otopk);
\draw[parrow, color=propgray!60] (otopk) -- (ogen);

\node[layerlabel=propred, above=0.3cm of oq, font=\small\bfseries] {\textbf{Original Hybrid (Superseded)}};

\begin{scope}[on background layer]
\node[boundbox=propgray, fit=(oq)(omf)(ovs)(orr)(otopk)(ogen), inner sep=8pt] {};
\end{scope}

% annotation for original hybrid
\draw[-{Stealth}, dashed, color=propred!70, thick]
  ([xshift=0.3cm]orr.east) -- ++(0.4,0)
  node[right, font=\scriptsize, color=propred, text width=1.8cm]
  {p95 = 25\,s\\weights dataset-\\dependent};

% ---- RIGHT: FINAL HYBRID (BM25+RRF) ----
\node[pbox=awsgray, minimum width=2.4cm, right=4.0cm of oq] (hq) {User Query};
\node[pbox=propred, below=of hq, minimum width=2.4cm] (hmf) {Metadata Filter\\(recency + source)};
\node[pbox=awsblue, below=of hmf, minimum width=2.4cm] (hvs) {Vector Similarity\\Search (Filtered)};
\node[pbox=awspurple, below=of hvs, minimum width=2.4cm] (hrr) {BM25 + RRF\\Fusion};
\node[pbox=awsgreen, below=of hrr, minimum width=2.4cm] (htopk) {Top-$k$ Selection};
\node[pbox=awsorange, below=of htopk, minimum width=2.4cm] (hgen) {LLM\\Generation};

\draw[parrow] (hq) -- (hmf);
\draw[parrow] (hmf) -- (hvs);
\draw[parrow] (hvs) -- (hrr);
\draw[parrow] (hrr) -- (htopk);
\draw[parrow] (htopk) -- (hgen);

\node[layerlabel=propblue, above=0.3cm of hq, font=\small\bfseries] {\textbf{Final Hybrid (BM25+RRF)}};

\begin{scope}[on background layer]
\node[boundbox=propblue, fit=(hq)(hmf)(hvs)(hrr)(htopk)(hgen), inner sep=8pt] {};
\end{scope}

% annotation for final hybrid
\draw[-{Stealth}, dashed, color=awspurple!70, thick]
  ([xshift=0.3cm]hrr.east) -- ++(0.4,0)
  node[right, font=\scriptsize, color=awspurple, text width=1.8cm]
  {p95 = 3.36\,s\\parameter-free};

\end{tikzpicture}%
}
\caption{Design evolution of retrieval strategies: baseline vector-only (left), original weighted hybrid with $\alpha$/$\beta$/$\gamma$ reranking (centre, superseded due to p95~=~25\,s tail latency), and final BM25+RRF hybrid pipeline (right).}
\label{fig:retrieval_comparison}
\end{figure}

\section{Hallucination Verification Strategy}
\label{sec:prompt}

TechPulse applies a three-layer verification approach to detect and flag potentially hallucinated outputs. The prompt construction and template details are presented in Section~\ref{sec:prompt_ch4}.

\begin{enumerate}[leftmargin=0.5in]
  \item \textbf{Prompt layer:} The system instruction enforces strict evidence grounding and abstention on insufficient context.

  \item \textbf{Evaluation layer:} RAGAS faithfulness scoring decomposes each generated response into atomic claims and verifies whether each claim is supported by at least one retrieved document.

  \item \textbf{Output layer --- Citation traceability:} At runtime, automated citation grounding verification analyzes each response sentence-by-sentence to compute a citation ratio---the fraction of sentences containing valid \texttt{[Source N]} references. Responses falling below the configurable citation grounding threshold are flagged as potentially hallucinated and replaced with an abstention message, with the original response preserved in metadata for review.
\end{enumerate}

\section{Optimization Objective}
\label{sec:optimization}

The retrieval task is framed as a constrained optimization problem:

\begin{equation}
  \max \; \text{Precision}(q) \quad \text{subject to} \quad \text{Latency} \leq \tau
  \label{eq:opt}
\end{equation}

where $\tau$ is the p95 latency target of 2 seconds. This formulation makes the precision--latency trade-off explicit and provides a principled basis for tuning retrieval depth and metadata filter selectivity.

\section{Tool Selection and Implementation Stack}
\label{sec:tools}

Table~\ref{tab:tools} summarizes the implemented tools and their justification. The stack has been adapted from the original proposal to operate within AWS Free Tier constraints, as detailed in Chapter~\ref{ch:progress}.

\begin{table}[H]
\caption{Implemented Tools, Services, and Justification.}
\label{tab:tools}
\centering
\begin{tabular}{p{3.2cm}p{3.2cm}p{6.6cm}}
\toprule
\textbf{Component} & \textbf{Tool / Service} & \textbf{Justification} \\
\midrule
Ingestion orchestration & AWS Lambda + EventBridge & Serverless, scheduled/event-driven, no cluster management \\
Message queue & Amazon SQS + DLQ & Decouples ingestion from embedding; failure resilience \\
Raw storage & Amazon S3 (medallion layout) & Cost-effective, durable, prefix-based lineage tracking \\
Embedding model & MiniLM (sentence-transformers) & Low latency, 384-dim, serverless-compatible \\
Vector + relational store & RDS PostgreSQL + pgvector & Unified hybrid retrieval; AWS Free Tier compatible \\
Generative model & Groq (Llama 3.1 8B Instant, primary); Bedrock (Amazon Nova Micro, fallback); Ollama / HuggingFace (local fallbacks) & Auto-cascade fallback chain; Groq free-tier API available across all environments \\
Frontend & React (S3 static hosting) & Lightweight, decoupled from backend \\
API routing & Amazon API Gateway (HTTP) & Managed, HTTPS, CORS support \\
CI/CD & GitHub Actions + SAM & Reproducible IaC deployments, automated testing \\
Monitoring & CloudWatch + SNS alarms & Observability, error alerts, health checks \\
Evaluation & RAGAS + statistical testing & Standardized RAG quality metrics + paired comparison \\
\bottomrule
\end{tabular}
\end{table}

\section{System Context Overview}
\label{sec:system_context}

Figure~\ref{fig:system_context} presents a high-level context view of TechPulse, showing external actors, data flows, and system boundaries before the detailed architecture is described in Chapter~\ref{ch:arch}.

% ---- TikZ: System Context Diagram ----
\begin{figure}[H]
\centering
\begin{tikzpicture}[node distance=1.0cm and 1.6cm, font=\small]

% Central system box
\node[rectangle, draw=awsblue, fill=awslightblue, rounded corners=6pt,
      minimum width=5.5cm, minimum height=3.2cm,
      align=center, font=\small\bfseries, line width=1.2pt]
  (sys) at (0,0) {\textbf{TechPulse}\\
  \scriptsize Hybrid RAG System\\
  \scriptsize (AWS Cloud + Docker Local)};

% External data sources (left) — pushed further out
\node[rectangle, draw=awsgreen!70, fill=awsgreen!15, rounded corners=4pt,
      text width=2.8cm, align=center, minimum height=0.65cm, font=\scriptsize]
  (arxiv) at (-6.2, 1.6) {ArXiv / Hacker News / DEV.to / GitHub / RSS};

% User (right top) — pushed further out
\node[rectangle, draw=awsorange!70, fill=awsorange!15, rounded corners=4pt,
      text width=2.6cm, align=center, minimum height=0.65cm, font=\scriptsize]
  (user) at (6.2, 1.0) {User};

% Evaluation system (right bottom) — pushed further out
\node[rectangle, draw=awspurple!70, fill=awspurple!15, rounded corners=4pt,
      text width=2.6cm, align=center, minimum height=0.65cm, font=\scriptsize]
  (eval) at (6.2, -1.0) {Evaluation (RAGAS + Stats)};

% Cloud / ops (bottom) — pushed further down
\node[rectangle, draw=propgray!70, fill=propgray!15, rounded corners=4pt,
      text width=2.8cm, align=center, minimum height=0.65cm, font=\scriptsize]
  (ops) at (0, -3.2) {Operators (CloudWatch / SNS Alerts)};

% Arrows with white-filled labels for visibility
\draw[-{Stealth}, thick, color=awsgreen!70]
  (arxiv.east) -- node[fill=white, inner sep=2pt, font=\scriptsize\bfseries, text=awsgreen!80] {raw content} (sys.west |- arxiv);

\draw[{Stealth}-{Stealth}, thick, color=awsorange!70]
  (user.west) -- node[fill=white, inner sep=2pt, font=\scriptsize\bfseries, text=awsorange!80, above, pos=0.5] {NL query}
                  node[fill=white, inner sep=2pt, font=\scriptsize\bfseries, text=awsorange!80, below, pos=0.5] {cited response}
  (sys.east |- user);

\draw[-{Stealth}, dashed, thick, color=awspurple!70]
  (sys.east |- eval) -- node[fill=white, inner sep=2pt, font=\scriptsize\bfseries, text=awspurple!80] {eval metrics} (eval.west);

\draw[-{Stealth}, dashed, thick, color=propgray!70]
  (sys.south) -- node[fill=white, inner sep=2pt, font=\scriptsize\bfseries, text=propgray!80, right] {drift alerts / health} (ops.north);

\end{tikzpicture}
\caption{System context diagram for TechPulse. Solid arrows represent data and query flows; dashed arrows represent monitoring and evaluation signals.}
\label{fig:system_context}
\end{figure}

\medskip
\noindent The following chapters translate these methodological decisions into a concrete system architecture (Chapter~4), evaluation framework (Chapter~5), implementation status (Chapter~6), and remaining work (Chapter~7).
